\input{estilo}
\begin{document}
\begin{figura}\begin{tikzpicture}[tic/.style={line width=1pt}]
  \foreach \x/\n in {0/1, 4/2} {
    \draw (\x,0) -- (\x,-6.2);
    \node[above] at (\x,0) {\var{P}\textsubscript{\n}};
  }
  % los caminos posibles: en línea llena el que se dibujó, en punteado la alternativa
  \draw[msg, suave, dashed] (0,-1.2) -- (4,-5.0);
  \draw[msg, suave, dashed] (4,-5.0) -- (0,-5.6);
  \draw[msg] (0,-1.2) -- node[cod, above, sloped] {get\_time} (4,-2.0);
  \draw[msg] (4,-4.0) -- node[cod, below, sloped] {ok} (0,-5.6);
  % marcas de tiempo
  \fill (0,-1.2) circle (2pt);
  \node[etiq, left=3pt] at (0,-1.2) {\var{T}\textsubscript{0}};
  \draw[tic] (-0.12,-3.4) -- (0.12,-3.4) node[etiq, right=2pt] {\var{T}\textsubscript{1}};
  \node[etiq, left=3pt] at (0,-5.6) {\var{T}\textsubscript{3}};
  \draw[tic] (3.88,-3.4) -- (4.12,-3.4);
  \draw[tic] (3.88,-5.0) -- (4.12,-5.0);
  \draw[llave, draw=segundo] (4.45,-2.0) -- (4.45,-5.0)
    node[midway, right=7pt, text=tinta] {\var{T}\textsubscript{2}};
\end{tikzpicture}\end{figura}
\end{document}
